\subsection*{一步一步建立直观}

\paragraph{1. 从速率记录重建总量}

如果只能读取瞬时流速，怎样恢复一段时间内的流入量？先用每个小区间起点的流速乘以持续时间，再把各段相加。右图的矩形把这个过程画出来。积分把“以离散记录估计累计”的方法推广为可控的极限。



\paragraph{2. 让分割变细，但不要混淆两种变量}

增加矩形数时，终点不变，改变的是估计精度；移动终点时，问题本身的累计范围改变。前者通向定积分的定义，后者产生面积函数。把这两种操作分开，就能理解为什么积分既是一个数，也可以随上限变化成为函数。



\paragraph{3. 由新增薄带得到基本定理}

将终点略向右移，新增面积约为端点函数值乘以移动距离。连续性保证这条窄带内部的高度不会突然偏离端点高度，于是面积函数的导数恢复原函数。再用任意原函数之差计算定积分，正是把这个局部事实连接到整个区间。



\subsection*{定义、命题与证明}

\paragraph{定义：可变上限的面积函数}

设 $f$ 在区间上连续，固定左端点 $a$，定义 $F(x)=\int_a^x f(t)\,dt$。这是随终点变化的有向累积量。

\paragraph{命题：连续速率的累积恢复原速率}

在区间内部，面积函数满足 $F'(x)=f(x)$。

\begin{proof}

把增量写成 $F(x+h)-F(x)=\int_x^{x+h}f(t)\,dt$。减去 $f(x)h$ 后，剩余量为 $\int_x^{x+h}(f(t)-f(x))\,dt$。

连续性保证足够短的区间内 $|f(t)-f(x)|$ 任意小。因此剩余量的绝对值至多为该小量乘 $|h|$。除以 $|h|$ 后取极限，差商与 $f(x)$ 的差趋零；正负步长都适用。

\end{proof}
